\section{Dense range and a bounded completed-zeta factor for the outside-plateau map}
\label{sec:dense-outside-transport}

\paragraph{Status and exact scope.}
The complete post-writing author reading and three independent
whole-proof readings have passed.
The seed in MA-PI-012 can be chosen so that its one fixed
operator has dense range in every stated completed theta ideal.
The proof also gives explicit preimages of Gaussian-regularized
smooth theta vectors and a bounded factorization of the completed
zeta convolution.  Neither construction is a bounded inverse on
the whole ideal, and the latter factor does not return the
original coordinate-flat local source.

Retain the exact one-place rational chart, actual ideal, and
Hilbert norms of (MA-PI.250)--(MA-PI.251).  In particular
\(\mathcal Mf(z)=\int f(x)e^{zx}\,dx\),
\(Vf(x)=2e^{x/2}f(e^x)\), and
\(J_{\mathscr H}=\overline{\mathscr I}^{\mathscr H}\).
Write \(D=d/dx\).  All parameters \(s\ge0\),
\(|b|<1/2\), and \(0<\alpha<1/2\) below are fixed when a
norm estimate is taken; no uniform estimate at the endpoint
widths is asserted.

The centered pole-removed completion and its actual kernel are
\[
 \begin{gathered}
 \mathfrak s=z+\tfrac12,\qquad
 \Psi(z)=\mathfrak s(\mathfrak s-1)
       \Gamma_{\mathbb R}(\mathfrak s)\zeta(\mathfrak s),\\
 \psi(x)=\frac1{2\pi}\int_{\mathbb R}\Psi(it)e^{-itx}\,dt
       \in\mathcal E,\qquad
 \mathscr I=\overline{\psi*\mathcal E}^{\,\mathcal E}.
 \end{gathered}
 \tag{MA-PI.280}
\]
The last equality is the actual source/ideal interface
MT-M18.4 and MT-M24.7--9 after centering, not a new
zero-location assumption.

\paragraph{MA-PI-013 (dense range with exact inverse limitations).}
There is a real nonnegative seed
\[
 h\in C_c^\infty((-\log2-2,-\log2-1)),\qquad
 H(z)=\mathcal Mh(z),
 \tag{MA-PI.281}
\]
whose Mellin transform has no zeros in \(|\Re z|<1/2\).
With the exact MA-PI-012 choice
\[
 \begin{gathered}
 g=(D+\tfrac12)h,\qquad G(z)=(\tfrac12-z)H(z),\\
 \gamma(z)=
  \frac{\Gamma_{\mathbb R}(1/2-z)}
       {\Gamma_{\mathbb R}(1/2+z)},\qquad
 r(z)=\tfrac12\gamma(z)\zeta(\tfrac12+z)G(z),
 \end{gathered}
 \tag{MA-PI.282}
\]
the same bounded injective operator \(R=C_k\), with
\(\mathcal Mk=r\), has
\[
             \overline{R\mathscr H}^{\,\mathscr H}
                    =J_{\mathscr H}
 \tag{MA-PI.283}
\]
for every space in (MA-PI.251).  All the outside-plateau
relations and full-equation failures of MA-PI-012 remain valid.

For each genuine smooth theta vector \(a_0=\Theta f\),
\(f\in\mathscr X_0\), there are explicitly defined vectors
\(u_n\) belonging to every one of these Hilbert spaces such
that
\[
 Ru_n=\gamma_n*a_0\longrightarrow a_0
       \quad\hbox{in each fixed }\mathscr H,\qquad
 \gamma_n(x)=\frac{n}{2\sqrt\pi}e^{-n^2x^2/4}.
 \tag{MA-PI.284}
\]
There is also one bounded convolution operator
\(B:\mathscr H\to\mathscr H\) in every such metric satisfying
\[
                    RB=C_\psi .
 \tag{MA-PI.285}
\]
Its exact multiplier is
\[
 \mathfrak b(z)=
 -\,\frac{2(z+\tfrac12)\Gamma_{\mathbb R}(z+\tfrac12)^2}
          {\Gamma_{\mathbb R}(\tfrac12-z)H(z)}.
 \tag{MA-PI.286}
\]
The range of \(R\) is proper and nonclosed in
\(J_{\mathscr H}\), and its inverse on that range is
unbounded, even when tested on original annular polar-free
local source vectors.  Finally, \(B\) is only a bounded
Hilbert-space factor: for an explicit such source vector
\(u\), \(V^{-1}Bu\) has a nonzero limit at the additive
origin.  Thus (MA-PI.285) is not a return to
\(\mathcal S_{00}\) or to \(\mathscr X_0\).

\begin{proof}
\emph{1. A compact smooth seed with a controlled zero set.}
For \(j\ge1\), put \(a_j=2^{-j}\) and let \(\nu_j\)
be the uniform probability measure on
\([-a_j/2,a_j/2]\).  The partial convolutions are supported
in \([-1/2,1/2]\).  For a continuous function on this
compact interval, adjoining the factors with \(j>N\)
changes its integral by at most its modulus of continuity
at \(\sum_{j>N}a_j/2\).  Their integrals therefore converge
to a positive functional of mass one, hence to a compactly
supported probability measure \(\nu\).  Its entire Mellin
transform is
\[
 U(z)=\prod_{j=1}^{\infty}
           \frac{\sinh(a_jz/2)}{a_jz/2}.
 \tag{MA-PI.287}
\]
The product converges uniformly on compact sets because
each factor is \(1+O(a_j^2|z|^2)\) there and
\(\sum_j a_j^2<\infty\).  This also proves the transform
identity by the preceding convergence of integrals.

On the imaginary axis each factor has modulus at most
\(\min(1,2^{j+1}/|t|)\).  For large \(|t|\), multiply
this estimate for
\(1\le j\le\lfloor\log_2|t|\rfloor-2\); the logarithm
of the resulting bound is a negative quadratic in
\(\log|t|\), up to a linear error.  Consequently there
are \(c_1,C_1>0\) such that
\[
 |U(it)|\le C_1
       \exp\{-c_1(\log(2+|t|))^2\}.
 \tag{MA-PI.288}
\]
In particular \(t^N U(it)\) is integrable for every
integer \(N\ge0\).  Fourier inversion shows that
\(\nu=u(x)\,dx\) with \(u\in C^\infty(\mathbb R)\).
Since \(\nu\) has compact support, \(u\in C_c^\infty\);
it is real, nonnegative, even, and has integral one.
Every derivative at its support endpoints is zero.

All zeros of an individual factor in (MA-PI.287) are
nonzero imaginary numbers.  If no factor vanishes at a
given point, the tail product is nonzero because the
sum of its deviations from one converges absolutely.
Thus \(U\) has no zeros off the imaginary axis.

Fix \(c>1/2\), put \(x_0=-\log2-3/2\), and set
\[
 h(x)=2e^{cx}u(2(x-x_0)),\qquad
 H(z)=e^{(z+c)x_0}U((z+c)/2).
 \tag{MA-PI.289}
\]
Its support is contained in
\([x_0-1/4,x_0+1/4]\), strictly inside the interval
in (MA-PI.281).  Its nonzero Mellin zeros lie on
\(\Re z=-c\), proving the required zero-free strip.
The rescaling by two ensures the strict support condition
of the original MA-PI-012 construction.

\emph{2. A uniform reciprocal estimate on every inner strip.}
In fact, if \(I\) is any compact real interval contained
in \((-c,\infty)\), there is \(C_I>0\) such that
\[
 |H(\sigma+it)|^{-1}
       \le \exp\{C_I(1+\log^2(2+|t|))\},
       \qquad \sigma\in I.
 \tag{MA-PI.290}
\]
To prove this, write \(w=\sigma+c+it\), so that
\(d=\Re w\) is bounded above and bounded below by
a positive number.  A factor of \(U(w/2)\) satisfies
\[
 \left|\frac{\sinh(a_jw/4)}{a_jw/4}\right|
       \ge\frac{d}{|w|},
 \tag{MA-PI.291}
\]
because
\(|\sinh(x+iy)|^2=\sinh^2x+\sin^2y\ge x^2\).
Choose \(J\) with \(J=O_I(\log(2+|t|))\) so that
\(|a_jw/4|\) is smaller than a fixed sufficiently small
constant for \(j>J\).  Multiplying (MA-PI.291) for
\(j\le J\) gives a lower bound
\(\exp\{-C_I(1+\log^2(2+|t|))\}\).
For the tail, the analytic logarithm near zero gives
\(\log|\sinh v/v|\ge-C|v|^2\).  Moreover
\(\sum_{j>J}|a_jw|^2\) is uniformly bounded.
The entire tail is therefore bounded below by a
positive constant.  The modulus of the exponential
factor in (MA-PI.289) is bounded above and below on
the chosen real interval.  This proves (MA-PI.290).
Cauchy's formula on fixed small disks in
\(\Re z>-c\) gives the same form of bound for every
fixed derivative of \(1/H\), using a slightly larger
compact interval.

\emph{3. Genuine source vectors and Gaussian preimages.}
For \(f\in\mathscr X_0\), the literal theta identity is
\[
 \mathcal M(\Theta f)(z)
       =\zeta(\tfrac12+z)\mathcal M(Vf)(z).
 \tag{MA-PI.292}
\]
It first follows in the absolutely convergent right
half-plane.  Here \(Vf\) is rapidly decreasing at
\(+\infty\) after every exponential weight.  At
\(-\infty\) it is \(O(e^{x/2})\), with the corresponding
derivative estimates, even without using \(f(0)=0\).
Thus \(\mathcal M(Vf)\) is holomorphic for
\(\Re z>-1/2\), and it and all its derivatives are
rapidly decreasing vertically on every fixed inner
substrip.  The pole at \(z=1/2\) is canceled by
\(\mathcal M(Vf)(1/2)=\int_{\mathbb R}f=0\).
Analytic continuation proves (MA-PI.292) on the
strip used below.  No coordinate-flatness of \(f\)
is needed.

Define, without division by zeta,
\[
 \begin{aligned}
 Q_f(z)&=
   \frac{2\mathcal M(Vf)(z)}
        {\gamma(z)(1/2-z)H(z)},\\
 Q_{f,n}(z)&=Q_f(z)e^{z^2/n^2},
                   \qquad |\Re z|<1/2.
 \end{aligned}
 \tag{MA-PI.293}
\]
All denominator factors are nonzero on this open
strip.  On every fixed closed inner substrip the
gamma quotient and its inverse have polynomial
growth, by Stirling's formula and compactness at
bounded height.  Equations (MA-PI.290) and
(MA-PI.293) consequently show that \(Q_{f,n}\),
with every fixed derivative, decreases faster than
every power of height, uniformly on that substrip.
Indeed \(e^{z^2/n^2}\) has modulus
\(e^{((\Re z)^2-(\Im z)^2)/n^2}\).

Let \(u_n\) be the inverse Mellin integral on
\(\Re z=0\).  Shift the contour to any two fixed
lines \(\pm\sigma\), \(0<\sigma<1/2\).
The horizontal edges tend to zero by the preceding
uniform estimates.  The same argument after
differentiating or multiplying the inverse integral
by powers of \(x\) gives
\[
                  e^{d x}u_n(x)\in\mathcal S(\mathbb R)
                  \quad\hbox{for every }|d|<1/2.
 \tag{MA-PI.294}
\]
For the polynomial factors one can equivalently
integrate by parts in the ordinate; all derivatives
of \(Q_{f,n}\) are rapidly decreasing on that line.
Choosing \(\sigma>|b|\) or \(\sigma>\alpha\)
proves \(u_n\in\mathscr H\) for every stated metric.
This assertion is not membership in \(\mathcal E\).

Multiplying the explicit expressions (MA-PI.282)
and (MA-PI.293) gives
\[
           r(z)Q_{f,n}(z)
             =\mathcal M(\Theta f)(z)e^{z^2/n^2}.
 \tag{MA-PI.295}
\]
The cancellation of zeta here is algebraic; it
remains valid at every zeta zero.  Ordinary Fourier
uniqueness on the central line, and then the actual
weighted convolution bounds, give the equality in
(MA-PI.284).  The Gaussian \(\gamma_n\) has integral
one and transform \(e^{z^2/n^2}\).

For completeness, it is an approximate identity in
each of the actual norms, not just on their smooth
quotients.  Translation is strongly continuous on
\(\mathscr H\), first on compactly supported smooth
functions and then by density.  The elementary
weight inequality gives
\[
 \|\tau_y v\|_{\mathscr H}
       \le C_s e^{\beta|y|}(1+y^2)^{s/4}
                    \|v\|_{\mathscr H},
 \quad
 \beta=|b|\ \hbox{or}\ \alpha.
 \tag{MA-PI.296}
\]
The integral of \(\gamma_n(y)\) times this bound
outside any fixed neighborhood of zero tends to
zero.  On that neighborhood strong continuity
controls \(\tau_yv-v\).  The integral inequality
for \(\gamma_n*v-v\) therefore proves (MA-PI.284).

The source interface used here reaches the full
ideal.  MT-M24.7--9 identifies its literal smooth
adelic image with the closed smooth ideal.
In the principal rational sector, average the
finite source over \(\widehat{\mathbb Z}^{\times}\).
It becomes a finite sum
\(\sum_\nu f_\nu\otimes\mathbf1_{q_\nu\widehat{\mathbb Z}}\),
\(q_\nu>0\).  Absorbing each \(q_\nu\) into
\(f_\nu(q_\nu t)\), then taking the even part,
gives one real Schwartz source.  The value at zero
is unchanged, and its integral is multiplied
termwise by \(q_\nu^{-1}\), exactly the finite
ball volume.  Both polar conditions, including
cancellations between summands, are preserved.
This is the actual MA-TR-101 argument.  Hence
every \(a_0\in\mathscr I\) is a vector \(\Theta f\)
of the form just treated.  Equation (MA-PI.284)
and the already proved \(R\mathscr H\subset
J_{\mathscr H}\) imply (MA-PI.283).

\emph{4. The bounded completed-zeta factor.}
The expression (MA-PI.286) is holomorphic in
\(|\Re z|<1/2\): the numerator gamma arguments
have positive real part there, the reciprocal
gamma is entire, and \(H\) is nonzero.
Stirling's formula gives on every closed inner
substrip, for suitable constants,
\[
 |\mathfrak b^{(j)}(\sigma+it)|
 \le C_{I,j}(2+|t|)^{C_{I,j}}
       \exp\{-\pi|t|/4+
                         C_{I,j}\log^2(2+|t|)\}.
 \tag{MA-PI.297}
\]
For \(j>0\) this follows either by Cauchy's formula
on a slightly wider substrip or by differentiating
the factors.  Its inverse Mellin kernel \(\ell\)
thus satisfies \(e^{dx}\ell\in\mathcal S\) for
every \(|d|<1/2\), by the same complete contour
argument as in step 3.  In particular
\[
 \int_{\mathbb R}|\ell(y)|
           e^{\beta|y|}(1+y^2)^{s/4}\,dy<\infty
             \qquad(\beta<1/2).
 \tag{MA-PI.298}
\]
Weighted Young and (MA-PI.296) prove that
\(B=C_\ell\) is bounded on every stated space.
Direct multiplication gives the entire-strip
identity
\[
        r(z)\mathfrak b(z)
           =(z+\tfrac12)(z-\tfrac12)
             \Gamma_{\mathbb R}(z+\tfrac12)
                       \zeta(z+\tfrac12)=\Psi(z).
 \tag{MA-PI.299}
\]
In particular the sign and factor \(-2(z+1/2)\)
in (MA-PI.286) are fixed by the pole-removal
polynomial, not by a chosen normalization at the
zeros.  Fourier uniqueness for the integrable
kernels proves \(k*\ell=\psi\), and their actual
bounded extensions prove (MA-PI.285).

This also supplies a second verification of the
full completed range.  The inclusion
\(\mathcal E\hookrightarrow\mathscr H\) is continuous
and dense, and \(C_\psi\) is bounded there.
Consequently
\[
 \overline{C_\psi\mathscr H}^{\,\mathscr H}
   =\overline{C_\psi\mathcal E}^{\,\mathscr H}
   =\overline{\mathscr I}^{\,\mathscr H}
   =J_{\mathscr H}.
 \tag{MA-PI.300}
\]
The middle equality follows from (MA-PI.280)
and continuity of the inclusion, not from a
spectral sum or an identification of orthogonal
projectors in different norms.

\emph{5. The inverse is unbounded on a genuine source core.}
Choose nonzero \(u_0,\beta_0\in C_c^\infty(\mathbb R)\)
with \(\mathcal M\beta_0(1/2)\ne0\).  Define
\[
 v_n(x)=e^{inx}u_0(x)-c_n\beta_0(x),\qquad
 c_n=\frac{\mathcal Mu_0(1/2+in)}
            {\mathcal M\beta_0(1/2)}.
 \tag{MA-PI.301}
\]
Integration by parts shows \(c_n=O_N(n^{-N})\)
for every \(N\).  Each \(v_n\) is compactly
supported and \(\mathcal Mv_n(1/2)=0\).
Thus \(V^{-1}v_n\) is a genuine even annular
source with value and integral zero.  In every
fixed metric
\(\|v_n\|_{\mathscr H}\to\|u_0\|_{\mathscr H}>0\).

The kernel \(k\) and all its derivatives have
the exponential Schwartz estimates of MA-PI-012.
For \(n\ne0\), integrate by parts \(N\) times
in
\[
 R(e^{in\cdot}u_0)(x)
       =e^{inx}\int k(y)u_0(x-y)e^{-iny}\,dy .
 \tag{MA-PI.302}
\]
For each \(x\) the differentiated factor has
compact \(y\)-support, so there are no boundary
terms.  The absolute value of every resulting
term is bounded by \(n^{-N}\) times a fixed
binomial coefficient and
\(|k^{(j)}|*|u_0^{(N-j)}|\).
Equation (MA-PI.296) and weighted Young bound
its \(\mathscr H\)-norm by a constant depending
on \(N\) and the fixed metric.  The correction
in (MA-PI.301) is rapidly small as well.  Hence
\[
             \|Rv_n\|_{\mathscr H}=O_N(n^{-N}).
 \tag{MA-PI.303}
\]
Normalizing \(v_n\) proves that the inverse on
the range is not bounded in that norm.
A bounded injective map between Hilbert spaces
with closed range has a bounded inverse onto
its range.  Thus \(R\mathscr H\) is nonclosed;
combined with (MA-PI.283) it is a proper dense
subspace of \(J_{\mathscr H}\).  In particular
there is no bounded right inverse
\(J_{\mathscr H}\to\mathscr H\).

\emph{6. The bounded factor fails the original origin condition.}
The seed estimate (MA-PI.290) holds a little
past the left boundary \(\Re z=-1/2\).
Choose \(0<\delta<\min(1,c-1/2)\).  In the
closed strip \(-1/2-\delta\le\Re z\le0\),
\(\mathfrak b\) is meromorphic with exactly
one pole, at \(z=-1/2\).  Indeed the next
numerator gamma pole is at \(z=-5/2\), and
all seed zeros lie at \(\Re z=-c\).
Since \(\Gamma_{\mathbb R}(w)=2/w+O(1)\)
at \(w=0\) and \(\Gamma_{\mathbb R}(1)=1\),
\[
           \operatorname*{Res}_{z=-1/2}
                  \mathfrak b(z)=-\frac8{H(-1/2)}.
 \tag{MA-PI.304}
\]
Let \(\beta\) be a nonzero nonnegative compact
smooth function and choose
\[
 u=(D+\tfrac12)\beta,\qquad
 \mathcal Mu(1/2)=0,\qquad
 \mathcal Mu(-1/2)=\mathcal M\beta(-1/2)>0.
 \tag{MA-PI.305}
\]
Then \(V^{-1}u\) is again an original annular
polar-free source.  The estimates proving
(MA-PI.297) also hold on the two sides of this
wider strip away from the pole, with rapid
vertical decrease after multiplication by
\(\mathcal Mu\).  Shift the inverse contour
from \(0\) to \(-1/2-\delta\); the horizontal
edges tend to zero.  The crossed residue gives
\[
 \begin{aligned}
 Bu(x)&=-\frac{8\mathcal Mu(-1/2)}{H(-1/2)}
                   e^{x/2}
            +O(e^{(1/2+\delta)x})\quad(x\to-\infty),\\
 \lim_{t\downarrow0}(V^{-1}Bu)(t)
       &=-\frac{4\mathcal Mu(-1/2)}{H(-1/2)}\ne0.
 \end{aligned}
 \tag{MA-PI.306}
\]
The sign follows by writing the upward central
contour as the upward left contour plus the
residue; the factor \(1/2\) in the second line
is the exact inverse chart.  This proves the
claimed failure of the original origin
condition without asserting any stronger
regularity of \(V^{-1}Bu\) at that point.
\end{proof}

\paragraph{Source and geometric placement.}
The new seed, lower bound, Gaussian construction,
weighted estimates, and boundary residue were
proved here.  The inherited source interfaces
are MA-TR-101's complete finite-ball absorption
argument and the actual MT-M24.7--9 image
identification.  The pole-removal definition
in MT-M12 and the unconditional MT-M18.1--6
transform argument fix (MA-PI.280); the GRH
and simplicity hypotheses of other MT-M12
claims are not used.  The local real gamma
factor is the Tate factor already checked in
MA-PI-012, with the same \(1/2\) chart scalar.
No new external zero-counting or zero-location
input is required.

In the inherited Grothendieck/functional-analytic
comparison, these are actual bounded convolution
maps and actual closures of the smooth theta
ideal.  The completed divisor and its finite
local lengths are unchanged.  Dense range is
not a topological isomorphism, an exact
conjugacy of cutoff projectors, or a nuclear
trace comparison.  The bounded factor \(B\)
regularizes by the full completed divisor,
and (MA-PI.306) explicitly prevents its use
as an unproved return to the original flat
local source.  This refinement does not decide
the FT002 arithmetic deficiency or close a
previously open unrestricted trace-factorization
question.
